%! Comparing Studies and Choosing a Method — Unit 1, Lesson 1.7 — Slides (Beamer)
% PILOT SHAPE (2026-09-06), notes in the MAIN IDEAS / NOTES table shape (2026-09-12).
% GRADUAL RELEASE deck: title → targets → warm-up → hook (unresolved) → I DO divider → two or
% three frames per notes ROW, labelled by row name and PROBLEM NUMBERS → WE DO Guided Practice
% (a live surface, un-answered) → spoken debrief with the cold check → close & START THE
% HOMEWORK, which is the individual practice. No group activity, no practice launch, no exit
% ticket. There is no spoiler rule: the notes frames ARE the direct instruction.
\documentclass[aspectratio=169, 11pt]{beamer}
\usepackage{saar-beamer}   % provides \ceruleanheader{} and \sectionlabel[color]{}

\newcommand{\redink}[1]{{\color{redacc}\bfseries #1}}

\begin{document}

% TITLE SLIDE (CourseName is not defined in beamer, so the name is literal)
{
\setbeamercolor{background canvas}{bg=cerulean}
\begin{frame}[plain]
  \vspace{2.8cm}\hspace{0.55cm}%
  \begin{minipage}{0.9\textwidth}
    {\color{goldacc}\bfseries\small LESSON 1.7}\\[0.5cm]
    {\color{white}\bfseries\LARGE Comparing Studies and Choosing a Method}\\[1.2cm]
    {\color{frostmid}\small Statistical Analysis \& Algebraic Reasoning~$\cdot$~Unit 1}
  \end{minipage}
\end{frame}
}

% ── LEARNING TARGETS ─────────────────────────────────────────────────────────
\begin{frame}[t, plain]
  \ceruleanheader{Today's targets}
  \sectionlabel{Lesson 1.7}
  By the end of today, I can\ldots
  \begin{itemize}\setlength{\itemsep}{5pt}
    \item compare an \textbf{observational study} with a \textbf{controlled experiment}, and
          write the strongest sentence each design earns.
    \item decide whether an experiment can be run at all --- or name it \textbf{unethical},
          \textbf{impossible}, or \textbf{impractical}.
    \item plan an experiment through the four stages of the \textbf{statistical cycle}.
    \item choose the \textbf{collection method} a question needs --- and keep the \emph{method}
          separate from the \emph{design}.
  \end{itemize}
  \begin{block}{How today works}
    Warm-up $\to$ \textbf{notes together, row by row} $\to$ \textbf{one pool study we work as a
    class} $\to$ debrief $\to$ \textbf{you start the homework here, alone}.
  \end{block}
\end{frame}

% ── WARM-UP ──────────────────────────────────────────────────────────────────
\begin{frame}[t, plain]
  \ceruleanheader{Warm-Up --- 5 minutes, silent}
  \sectionlabel{back to Riverbend's Study Center}
  \vspace{0.2cm}
  \begin{center}
    \begin{minipage}{0.88\textwidth}
      \textbf{Last quarter she watched:} 150 records --- 60 used the Study Center (39 passed
      every class), 90 did not (72 passed). \textbf{This quarter she assigned:} a generator
      split 120 volunteers 60/60 --- 75\% against 55\%.
      \begin{enumerate}\setlength{\itemsep}{6pt}
        \item What percent of the 60 users passed every class?
        \item What percent of the 90 non-users passed?
        \item Both studies side by side: the two gaps, the group sizes --- and who put each
              student in a group \emph{last} quarter?
      \end{enumerate}
    \end{minipage}
  \end{center}
  \vspace{0.3cm}
  \begin{center}
    {\color{cerulean}\bfseries 60 and 90 \quad versus \quad 60 and 60. That stays on the board
    all period.}
  \end{center}
\end{frame}

% ── HOOK — leave it UNRESOLVED; problem 15 (row 4, The Method) settles it ────
{
\setbeamercolor{background canvas}{bg=cerulean}
\begin{frame}[plain]
  \vspace{1.2cm}
  \begin{center}
    {\color{frostmid}\bfseries\small A GENERATOR SPLIT 120 VOLUNTEERS 60/60. THEN SHE SAT IN THE
    STUDY CENTER AND WATCHED.}\\[0.8cm]
    {\color{white}\bfseries\Large ``She just watched. So this quarter's study was an
    observational study, the same as last quarter's.''}\\[0.9cm]
    {\color{goldacc}\bfseries\large Is the classmate right?}\\[0.5cm]
    {\color{frostmid}\small Vote: agree or disagree. We settle this at problem 15.}
  \end{center}
\end{frame}
}

% ── I DO — direct instruction begins ─────────────────────────────────────────
{
\setbeamercolor{background canvas}{bg=cerulean}
\begin{frame}[plain]
  \vspace{1.8cm}\hspace{0.8cm}%
  \begin{minipage}{0.86\textwidth}
    {\color{goldacc}\bfseries\small GUIDED NOTES \ \textbullet\ \ I DO}\\[0.7cm]
    {\color{white}\bfseries\Large Open to your Guided Notes page. We work the table row by
    row --- fill each term as we name it, not before.}\\[0.9cm]
    {\color{frostmid}\small Five terms go in the vocabulary box today. Write each one the
    moment we use it.}
  \end{minipage}
\end{frame}
}

% ── ROW 1a. TWO DESIGNS ──────────────────────────────────────────────────────
\begin{frame}[t, plain]
  \ceruleanheader{Two designs, one question}
  \sectionlabel{notes row 1 $\cdot$ Two Designs $\cdot$ problems 1--3}
  \vspace{4pt}
  \begin{columns}[t]
    \begin{column}{0.47\textwidth}
      \begin{block}{Observational study}
        Measures or asks \textbf{without assigning} anyone to a group --- the groups already
        exist.
      \end{block}
      \vspace{4pt}
      {\small Riverbend, last quarter: 60 used it, 90 did not. \textbf{They} sorted themselves.}
    \end{column}
    \begin{column}{0.47\textwidth}
      \begin{block}{Controlled experiment}
        \textbf{Assigns} the groups by chance and imposes a treatment on at least one of them.
      \end{block}
      \vspace{4pt}
      {\small Riverbend, this quarter: a generator split 120 into 60/60.}
    \end{column}
  \end{columns}
  \vspace{12pt}
  \begin{center}
    {\large One question sorts them: \redink{who decided which group each individual is in?}}\\[6pt]
    {\small\color{slate} Every other row of the table --- treatment, groups alike, the verb ---
    follows from that answer.}
  \end{center}
\end{frame}

% ── ROW 1b. THE VERB EACH DESIGN EARNS ───────────────────────────────────────
\begin{frame}[t, plain]
  \ceruleanheader{The design decides the verb}
  \sectionlabel{notes row 1 $\cdot$ problems 2--5 $\cdot$ two studies, opposite results}
  \vspace{2pt}
  \begin{center}
    \small
    \renewcommand{\arraystretch}{1.3}
    \begin{tabular}{@{} l l l @{}}
    \hline
    \textbf{The study} & \textbf{What it found} & \textbf{What it may say} \\
    \hline
    Watched 150 records & 65\% of users vs.\ 80\% of non-users & \emph{is associated with} \\
    Assigned 120 volunteers & 75\% treated vs.\ 55\% control & \emph{caused} \\
    \hline
    \end{tabular}
  \end{center}
  \vspace{6pt}
  \begin{itemize}\setlength{\itemsep}{4pt}
    \item Same school, same Study Center, \textbf{opposite conclusions}. Which one do you
          believe? \redink{This quarter's.}
    \item Chance built two groups that started out alike, so the treatment is the only
          difference left to explain the gap.
    \item Problem 5: 60 and 90 against 60 and 60. \redink{The sizes give away who built the
          groups.}
  \end{itemize}
\end{frame}

% ── ROW 2. WHEN YOU CANNOT RUN AN EXPERIMENT (the trap) ──────────────────────
\begin{frame}[t, plain]
  \ceruleanheader{When you cannot run an experiment}
  \sectionlabel{notes row 2 $\cdot$ Experiment $\cdot$ problems 6--9 $\cdot$ problem 8 is the trap}
  \vspace{2pt}
  \begin{center}
    \small
    \renewcommand{\arraystretch}{1.3}
    \begin{tabular}{@{} l l @{}}
    \hline
    \textbf{Reason} & \textbf{What it means} \\
    \hline
    Unethical   & Assigning the treatment would harm someone. \\
    Impossible  & The treatment is something a person already \emph{is} or \emph{has}. \\
    Impractical & Too costly, too slow, or nobody would agree to be assigned. \\
    \hline
    \end{tabular}
  \end{center}
  \vspace{6pt}
  \begin{block}{Problem 8 --- vote first: the 45-minute bus ride}
    A classmate says \emph{``unethical.''} Who is harmed by studying it? Nobody. You simply
    \redink{cannot assign where a family lives} --- the reason is \redink{impossible}.
  \end{block}
  \vspace{4pt}
  \begin{center}
    {\color{cerulean}\bfseries ``We could not run an experiment'' is the reason for a weaker
    verb --- never an excuse for a stronger one.}
  \end{center}
\end{frame}

% ── ROW 3. THE CYCLE — fast ──────────────────────────────────────────────────
\begin{frame}[t, plain]
  \ceruleanheader{Planning an experiment: the same four stages}
  \sectionlabel{notes row 3 $\cdot$ The Cycle $\cdot$ problems 10--11}
  \vspace{4pt}
  \begin{center}
    \small
    \renewcommand{\arraystretch}{1.3}
    \begin{tabular}{@{} l l @{}}
    \hline
    \textbf{Stage} & \textbf{The weekly check-in study} \\
    \hline
    Formulate         & Does a weekly check-in raise assignments turned in? \\
    \textbf{Collect}  & \textbf{The units are the volunteers; a generator \emph{assigns} them} \\
    Represent         & A table: each group's count and percent \\
    Analyze \& report & Compare the two percents; use the verb the design earns \\
    \hline
    \end{tabular}
  \end{center}
  \vspace{8pt}
  \begin{center}
    {\large Lesson 1.5 \emph{sampled} the units. Today we \redink{assign} them.}\\[4pt]
    {\small\color{slate} One word changes. The other three stages are the ones you already
    know.}
  \end{center}
\end{frame}

% ── ROW 4a. FIVE WAYS TO COLLECT ─────────────────────────────────────────────
\begin{frame}[t, plain]
  \ceruleanheader{Five ways to collect the data}
  \sectionlabel{notes row 4 $\cdot$ The Method $\cdot$ problems 12--14}
  \vspace{2pt}
  \begin{center}
    \small
    \renewcommand{\arraystretch}{1.3}
    \begin{tabular}{@{} l l @{}}
    \hline
    \textbf{Method} & \textbf{Best when you need\ldots} \\
    \hline
    Survey           & the same fixed questions from many people \\
    Interview        & depth from one person --- the reasons, the \emph{why} \\
    Focus group      & a small group building on each other's answers \\
    Observation      & what people actually \emph{do}, not what they say they do \\
    Content analysis & records that already exist --- change over time \\
    \hline
    \end{tabular}
  \end{center}
  \vspace{8pt}
  \begin{center}
    {\color{cerulean}\bfseries The question picks the method. The method never picks the
    design.}
  \end{center}
\end{frame}

% ── ROW 4b. THE CRUX ─────────────────────────────────────────────────────────
\begin{frame}[t, plain]
  \ceruleanheader{The method is not the design}
  \sectionlabel[redacc]{notes row 4 $\cdot$ problem 15 $\cdot$ the whole point of today}
  \vspace{4pt}
  \begin{columns}[t]
    \begin{column}{0.47\textwidth}
      \begin{block}{Who built the groups?}
        A random number generator, 60/60. \\[3pt]
        $\Rightarrow$ the \textbf{design} is a \redink{controlled experiment}.
      \end{block}
    \end{column}
    \begin{column}{0.47\textwidth}
      \begin{block}{How were the answers recorded?}
        She sat in the Study Center and watched. \\[3pt]
        $\Rightarrow$ the \textbf{method} is \redink{observation}.
      \end{block}
    \end{column}
  \end{columns}
  \vspace{10pt}
  \begin{itemize}\setlength{\itemsep}{4pt}
    \item \textbf{Problem 15 --- vote again: is the classmate right?}
    \item \emph{Observation} is a method. An \emph{observational study} is a design. They are
          not the same word twice.
  \end{itemize}
  \begin{block}{PS.DC.3e}
    \textbf{An experiment can collect its data by observation.} How the data arrived
    \redink{never} decides who built the groups.
  \end{block}
\end{frame}

% ── WE DO — a LIVE surface: task and data only, no answers ───────────────────
\begin{frame}[t, plain]
  \ceruleanheader{Guided Practice: Harbor Point Community Pool}
  \sectionlabel[goldacc]{We do $\cdot$ problems 16--19 $\cdot$ you hold the pen}
  1{,}500 members. Two studies already on the desk (1.5: surveyed 200; 1.6: assigned 90).
  \textbf{Now:} 120 volunteers, a generator sends 60 to a weekly water-aerobics class and 60 to
  no class for six weeks. \textbf{42} of the class group and \textbf{30} of the others report
  less joint pain.
  \begin{itemize}\setlength{\itemsep}{5pt}
    \item[16.] Label the two earlier studies by design --- and who built the groups in each?
    \item[17.] The two rates, the gap, and the strongest sentence this study earns.
    \item[18.] New question: \emph{do members who already have arthritis sleep worse?} Can she
          assign it? Which reason? What does she run, with which verb?
    \item[19.] She asked all 120 the same three questions; she also wants to know \emph{why}
          some quit. Name both methods --- does either change the design?
  \end{itemize}
  \begin{block}{The question behind the question}
    If 19 changed your answer to ``observational study,'' you are still at the hook vote.
  \end{block}
\end{frame}

% ── DEBRIEF ──────────────────────────────────────────────────────────────────
\begin{frame}[t, plain]
  \ceruleanheader{Debrief: who built the groups?}
  \sectionlabel{10 minutes $\cdot$ pens down, papers up --- we say these aloud}
  \vspace{4pt}
  \begin{enumerate}\setlength{\itemsep}{5pt}
    \item The counselor watched --- and a generator built her groups. Say the sentence.
    \item Problem 8: who could hand out a 45-minute bus ride? Then what is the reason?
    \item Problems 2 and 3: same Study Center, two verbs. What bought \emph{caused}?
  \end{enumerate}
  \vspace{6pt}
  \begin{block}{Cold check --- hands down, thirty seconds}
    Bayside Middle School. \textbf{Study A:} the counselor compared the 150 students in band
    with the 600 who are not --- 90 of the band students had a B average. \textbf{Study B:} 80
    sixth-graders volunteered; a generator sent 40 to a silent reading block and 40 to homeroom
    --- 70\% improved against 40\%. \textbf{Which study may say the program \emph{caused} the
    difference --- and what percent of the band students had a B average?}
  \end{block}
\end{frame}

% ── YOU DO — the homework is the individual practice, started in class ───────
\begin{frame}[t, plain]
  \ceruleanheader{You do: start the homework}
  \sectionlabel[goldacc]{Last 10 minutes $\cdot$ on your own $\cdot$ finish it at home}
  \begin{itemize}\setlength{\itemsep}{5pt}
    \item \textbf{Millbrook Public Library} --- 1{,}200 card holders, and now a \emph{third}
          study: 300 volunteers, 150/150, a weekly evening book club.
    \item Sort the two earlier studies; compute the book-club rates and the gap.
    \item Plan the study through the cycle; decide two new questions --- assign or not?
    \item Match all five collection methods to five jobs.
    \item Then the two claims: she \emph{watched} the book club --- and the 48-point gap.
  \end{itemize}
  \begin{block}{Silent and alone --- I am circulating. Packet pages, scored out of 12, due at the start of the first class after two study halls.}
    Start with \textbf{1, 2, and 6}. Item 6 is the one I am reading over your shoulder --- if
    you called the book-club study observational, flag me before you leave.
  \end{block}
\end{frame}

% ── CLOSE ────────────────────────────────────────────────────────────────────
{
\setbeamercolor{background canvas}{bg=cerulean}
\begin{frame}[plain]
  \vspace{1.5cm}\hspace{0.8cm}%
  \begin{minipage}{0.86\textwidth}
    {\color{goldacc}\bfseries\small BEFORE YOU GO}\\[0.7cm]
    {\color{white}\bfseries\Large Who built the groups decides the design. How the answers were
    recorded decides the method. The biggest gap in this unit is the one you should trust
    least.}\\[0.9cm]
    {\color{frostmid}\small Next: Lesson 1.8 --- \emph{You Run the Study}. Today you judged
    other people's designs. Next time the question is yours, and every decision on this page is
    one you have to make.}
  \end{minipage}
\end{frame}
}

\end{document}
